\chapter*{How to use this book}
\addcontentsline{toc}{chapter}{How to use this book}

Work at a computer and compile often. Reading a code sample is useful, but
changing it and examining the new PDF is where understanding becomes
practical. Keep a separate practice project so that experiments cannot damage
important work.

For a faster, practical route, compile the first complete code panel in each
chapter before reading every explanation. Make the changes named in the Code
Lab, compare the PDFs, and return to the surrounding text when the result
surprises you. This keeps the source file open while you learn.

\section*{Choose a working path}

The beginner-friendly path uses Overleaf in a web browser. It avoids an
installation during the first lesson and places source, compiler messages,
and PDF preview in one project. The local-computer path uses TeX Live on macOS
or Linux, or MiKTeX/TeX Live on Windows. TeXstudio and Visual Studio Code with
\LaTeX{} Workshop are optional editors; \pkg{latexmk} is the preferred automated
build tool. The \LaTeX{} source is the same on both paths.

Start with whichever path lowers the immediate barrier. Later, try the other
path so that you understand how to keep and compile your own complete source.

\section*{Use the spiral project}

At the end of each chapter, improve the continuing measurement-methods
article. Do not jump to the final version at once. The value lies in seeing
why each layer is added and which problem it solves.

\section*{Make errors on purpose}

Keep a working copy, make the requested error, compile it, and read the first
meaningful message. Then repair the source. A controlled error is a safe way
to learn what a message means before it appears in a deadline-driven project.

\section*{Follow external instructions when they govern}

A journal class, thesis manual, funder, or client may impose rules that differ
from the generic examples. Verify the current instructions and change only
what they require. The book explains the reasoning needed to make those
changes without dismantling the document's structure.
